% Edward B. Baker III, with OpenAI Codex assistance, 25 November 2026.
% Inputs: repository HEAD 5ff7014 and the dated neck-geometry derivations.
% Exact serving model and reasoning effort were unavailable.

\section{Positive Schouten representatives and boundary charge}
\label{alt:gauge}

Let \(M\) be a smooth closed oriented homotopy four-sphere. Write
\(E=E_+=E_-\) for its common closed chiral Weyl energy, and suppose
\(Y=Y([g])>0\) and \(E\le 8\pi^2-\delta\), \(\delta>0\). For a
representative \(h\) use the Schouten tensor \(P_h=(\operatorname{Ric}_h-R_hh/6)/2\),
and put \(p_h=\sigma_2(h^{-1}P_h)\) and \(T_h=\operatorname{Ric}_{0,h}\).
The elementary symmetric function and Chern--Gauss--Bonnet give
\begin{equation}\label{alt:schouten}
 p_h=\frac{R_h^2}{96}-\frac{|T_h|^2}{8},\qquad
 q:=\int_M p_h\,dV_h=\frac{8\pi^2-E}{2}\ge \frac{\delta}{2}.
\end{equation}
In the unit-volume Yamabe representative \(g_Y\),
\(R_{g_Y}=Y\) and \(Y^2=96q+12\int_M|T_{g_Y}|^2\,dV_{g_Y}\).
In particular \(\sqrt{48\delta}\le Y\le Y_{\mathrm{rd}}=8\sqrt6\pi\).

Theorem A and Corollary B of Chang--Gursky--Yang give a smooth
representative \(h\) with \(R_h>0\) and \(p_h>0\)
\cite[Theorem A, Corollary B and Lemma 1.1]{ChangGurskyYang2002PositiveRicci}.
Their Schouten tensor is \(2P_h\). In the same class one may choose
the unit-volume representative satisfying \(p_h=q\), by their
constant-\(\sigma_2\) existence theorem
\cite[Theorem 1.2 and equation (1.6)]{ChangGurskyYang2003}.
Positivity of \(p_h\) prevents \(R_h\) from vanishing by
\eqref{alt:schouten}. Connectedness and the positive Yamabe
constant, tested against the constant function, give \(R_h>0\).
Uniqueness after volume normalization applies in a fixed nonround
class, and the solutions in the round class are round
\cite[Theorem 1.5]{GurskyStreets2018TwoYamabe}.

For \(S=R_h/6>0\) and least eigenvalue \(a\) of \(h^{-1}P_h\),
Cauchy--Schwarz on the other three eigenvalues gives
\[
 p_h\le a(S-a)+(S-a)^2/3=(S-a)(S+2a)/3.
\]
Thus \(m=S+2a=\lambda_{\min}(\operatorname{Ric}_h)>0\), and
\(p_h\le m(3S-m)/6\le mS/2\). Also \(|P_h|^2<S^2\), because
\(S^2-|P_h|^2=2p_h\). These inequalities give
\begin{equation}\label{alt:riccone}
 R_h\ge\sqrt{96p_h},\qquad
 \lambda_{\min}(\operatorname{Ric}_h)\ge\frac{12p_h}{R_h},
 \qquad \operatorname{Ric}_h<\frac{R_h}{2}h.
\end{equation}
At fixed \(p_h=q>0\) the algebraic bound can degenerate.
The Ricci eigenvalues \(m,(R-m)/3,(R-m)/3,(R-m)/3\), with
\(m=(R-\sqrt{R^2-96q})/4\), have \(m\sim12q/R\). The modified
prescription
\(p_h-\alpha |W_h|_{\mathrm{tens}}^2/16=\lambda\) has
\(\lambda=q-\alpha E/2\ge\delta/4\) at volume one for
\(\alpha=\delta/(16\pi^2)\)
\cite[Theorem 1.1 and equation (1.4)]{ChangGurskyYang2003}.
The same algebraic degeneration is possible with a zero Weyl block.

\subsection{Physical metric estimates}

The chiral Weyl measures and \(|B|\,dV\) are conformally invariant
on each fixed domain. Physical balls and their energy quantile radii
depend on the representative. In any representative one has
\begin{equation}\label{alt:full-curvature}
 \begin{split}
 \int_M |T_h|^2\,dV_h&=\tfrac1{12}\int_M R_h^2\,dV_h-8q,\\
 \int_M |\operatorname{Ric}_h|^2\,dV_h
   &=\tfrac13\int_M R_h^2\,dV_h-8q,\\
 \int_M |\operatorname{Rm}_h|_{\mathrm{tens}}^2\,dV_h
   &=8E-16q+\tfrac13\int_M R_h^2\,dV_h .
 \end{split}
\end{equation}
The first identity follows from \eqref{alt:schouten}. The second uses
the trace decomposition, and the third uses
\(|\operatorname{Rm}|_{\mathrm{tens}}^2
=|W|_{\mathrm{tens}}^2+2|T|^2+R^2/6\).
In the Yamabe gauge \(\int R^2=Y^2\le384\pi^2\). In a positive-Ricci
representative with volume \(V\) and diameter \(D\), Bishop--Gromov gives
\[
 \frac{V}{D^4}r^4\le V_x(r):=\operatorname{Vol}_h B_h(x,r)
 \le \omega_4r^4\quad(0<r\le D),\qquad \omega_4=\pi^2/2.
\]
A scalar upper bound \(R_h\le R_*\) in the constant-\(p_h\) gauge
would give \(D\le\pi\sqrt{R_*/(4q)}\) by \eqref{alt:riccone} and
Bonnet--Myers.

\begin{proposition}[Scalar energy and physical ball bounds]
\label{alt:physical-balls}
Suppose \(h\) has unit volume, \(\operatorname{Ric}_h\ge0\),
\(Y([h])\ge y_0>0\) and \(\int_M R_h^2\,dV_h\le L_*^2\).
Then its diameter is bounded above in terms of \(y_0,L_*\), and
\(c(y_0,L_*)r^4\le V_x(r)\le\omega_4r^4\) for \(0<r\le D\).
For a complete connected noncompact four-manifold with
nonnegative Ricci curvature, finite scalar \(L^2\) energy
and the same Yamabe inequality on compactly supported tests,
the asymptotic volume ratio is at least \((y_0/Y_{\mathrm{rd}})^2\).
\end{proposition}

\begin{proof}
Use a cutoff supported on \(B(x,r)\), equal to one on \(B(x,r/2)\),
with gradient bound \(2/r\). Bishop doubling gives
\(V_x(r/2)\ge V_x(r)/16\). The Yamabe inequality and Hölder yield
\[
 \frac{y_0}{4}\le \frac{24}{r^2}V_x(r)^{1/2}
 +\left(\int_{B(x,r)}R_h^2\,dV_h\right)^{1/2}.
\]
Thus each ball has volume at least \(\kappa r^4\),
\(\kappa=(y_0/192)^2\), or scalar \(L^2\) norm greater than
\(y_0/8\). Pack disjoint unit balls centered two units apart along a
minimizing diameter segment. The two kinds of balls number at most
\(\kappa^{-1}+64L_*^2/y_0^2\) altogether, which bounds \(D\) and gives
the first assertion by Bishop--Gromov.

In the noncompact case let \(\alpha_\infty\) be the asymptotic
volume ratio. Balls of radius \(r\) centered at distance \(2r\) along
a ray have scalar energy tending to zero, which first gives
\(\alpha_\infty\ge\kappa/(81\omega_4)>0\).
Apply the compact-support Yamabe inequality to \(f(d(o,\cdot)/r)\)
for \(f\in C_c^\infty((0,\infty))\). The pushforward of
\(r^{-4}dV\) by \(d(o,\cdot)/r\) tends on compact subintervals to
\(\alpha_\infty\,d(\omega_4s^4)\), by convergence of the radial volume
distribution. The distance gradient has norm one almost everywhere.
The scalar term divided by \(r^2\) tends to zero by Hölder and
finiteness of the scalar energy on the escaping supports. The limit
inequality bounds \(y_0\) by \(\sqrt{\alpha_\infty}\) times the
Euclidean Yamabe quotient of the radial test. Cutoffs of the
stereographic round conformal factor have quotients tending to
\(Y_{\mathrm{rd}}\) \cite[\S3, Theorems 3.2--3.3]{LeeParker1987Yamabe}.
\end{proof}

For a supplied complete noncompact smooth pointed limit at
vanishing scales in these constant-\(p_h\) gauges, assume the
uniform scalar-energy bound of Proposition~\ref{alt:physical-balls}.
Nonnegative Ricci curvature and the compact-support Yamabe
inequality pass to the limit. The
scalar \(L^2\) budget passes on exhaustion and the constant \(p_h\)
rescales to zero. Its volume ratio is at least
\(\delta/(8\pi^2)\). A single smooth ALE end with link \(S^3\)
has volume ratio one, in which case Bishop--Gromov rigidity gives
Euclidean space.

Positivity alone leaves the scalar-energy and diameter estimates
open. For example, the rotational ellipsoids
\(z^2/L^2+|y|^2=1\) in \(\mathbb R^5\) are strictly convex and
conformally round. In the warped-product expression
\(ds^2+a(s)^2\gamma_{S^3}\), this follows by using
\(d\log\rho=ds/a(s)\), with smooth pole expansions
\(a(s)=s+O(s^3)\) in pole distance. If \(k_r,k_t\) are radial
and tangential sectional curvatures, then
\(p_h=3k_rk_t/2>0\). On \(|z|\le L/2\),
\(a=\sqrt{1-z^2/L^2}\) and \(k_t=(1-a_s^2)/a^2\) have uniform
positive lower bounds. The volume is comparable to \(L\) and
\(\int R^2\ge cL\). Unit-volume rescaling leaves that divergence
and gives diameter at least \(2L/V^{1/4}\). A unit ball in that
rescaling lies in an original \(z\)-slab of width \(2V^{1/4}\),
so has unit-volume measure at most \(CV^{-3/4}\). This family has
nonconstant \(p_h\). The constant-\(p_h\) choice in its round class
is round. For varying constant-\(p_h\) classes,
\eqref{alt:full-curvature} and Proposition~\ref{alt:physical-balls}
use a further uniform scalar-energy assumption. Transferring a
physical radius selection between gauges also requires distance
comparison or a new selection in the chosen metric.

\subsection{Conformal boundary current and its modular consequence}

Fix an actual compact smooth domain \(U\subset M\).
Let \(h=e^{2u}g\), \(g_t=e^{2tu}g\), and
\(N_t=(\operatorname{tr}_{g_t}P_{g_t})g_t-P_{g_t}\).
The contracted Bianchi identity gives \(\operatorname{div}N_t=0\).
The infinitesimal Schouten law
\(\partial_tP_{g_t}=-\nabla^2_{g_t}u\) and
\(\operatorname{tr}(T_1(K)K)=2\sigma_2(K)\) for
\(K=g_t^{-1}P_{g_t}\) give
\begin{equation}\label{alt:newton}
 \partial_t(p_{g_t}\,dV_{g_t})
 =-\operatorname{div}_{g_t}(N_t^{\sharp_{g_t}}\nabla^{g_t}u)\,dV_{g_t}.
\end{equation}
Here the term \(-4up_{g_t}\) from differentiating the inverse
metric cancels the volume derivative. These are the Newton tensor
identities of \cite[\S2.1, equations (2-1)--(2-5) and Lemma 2.1]{GurskyStreets2018TwoYamabe}
with the sign appropriate to \(g_t=e^{2tu}g\). Write
\[
 q_g(U)=\int_U p_g\,dV_g,\qquad
 J=\sum_{\Sigma\subset\partial U}
   \int_0^1\!\int_\Sigma N_t(\nu_t,\nabla^{g_t}u)\,dS_{g_t}\,dt,
 \qquad q_h(U)-q_g(U)=-J ,
\]
using the outward \(g_t\)-unit normals. In this integral the changes
of the normal, gradient and area contribute factors
\(e^{-tu},e^{-2tu},e^{3tu}\), whose product is one. Integrating
\(P_{g_t}=P_g-t\nabla_g^2 u+t^2du\otimes du
-t^2|du|_g^2g/2\) gives the exact face current
\begin{equation}\label{alt:face-current}
 J_\Sigma=\int_\Sigma\!
  \left[N_0(\nu,\nabla u)
   +\tfrac12\bigl(\nabla^2u(\nu,\nabla u)-(\Delta_g u)u_\nu\bigr)
   -\tfrac12|\nabla u|^2u_\nu\right]dS_g .
\end{equation}
Here \(\Delta_g=\nabla^i\nabla_i\). If \(f=u|_\Sigma\),
\(v=\nabla_\Sigma f\), \(a=u_\nu\), and
\(II(X,Y)=g(\nabla_X\nu,Y)\), the same integral has integrand
\[
 aN_0(\nu,\nu)+N_0(\nu,v)-a\Delta_\Sigma f
 -\tfrac12 II(v,v)-\tfrac12(\operatorname{tr}II)a^2
 -\tfrac12a(a^2+|v|^2).
\]
Indeed the normal-normal Hessian cancels and the mixed Hessian
contributes \(\langle\nabla_\Sigma a,v\rangle-II(v,v)\).
The two integrals of the tangential Laplacian combine by integration
by parts. Hölder bounds \(|J_\Sigma|\) by
\[
 C\bigl(A_\Sigma D_\Sigma+\|II\|_3 A_\Sigma^2+A_\Sigma^3\bigr),
 \quad A_\Sigma=\|\,|a|+|v|\,\|_3,\quad
 D_\Sigma=\|N_0(\nu,\cdot)\|_{3/2}+\|\Delta_\Sigma f\|_{3/2}.
\]
All the norms in this bound are invariant under constant rescaling.
For paired faces with a common forward transverse direction,
\eqref{alt:face-current} gives the difference of the two integrals
in that direction. Cancellation of the signed sum is possible even
when individual integrals are large.

For either chirality put
\(I_\pm(U,h)=E_\pm(h;U)+2q_h(U)\), and denote by
\(C_\pm(\partial U,h)\) the oriented sum of the Chern--Simons
integrals of the physical Levi-Civita half-spin connections in
boundary frames. Use the unnormalized anti-Hermitian fundamental
trace and set \(P_0=8\pi^2\), \(s_+=-1\), \(s_-=1\).
The homotopy sphere is spin because \(H^2(M;\mathbb Z/2)=0\).
Its half-spin bundles restrict trivially to the boundary
since \(\pi_j(\mathrm{SU}(2))=0\) for \(j=0,1,2\).
The diagonal Riemann curvature block is
\(W^\pm+R_h\operatorname{Id}/12\) and the off-diagonal block
has squared norm \(|T_h|^2/4\). The factor \(1/2\) in the
Clifford curvature map is paired with
\(-\operatorname{tr}(c(\omega)c(\eta))=4\langle\omega,\eta\rangle\)
for two-forms of the acting chirality
\cite[\S2.3]{Baker2026DiracYamabe}. Evaluating
\(s_\pm\operatorname{tr}(F_\pm\wedge F_\pm)\) gives the difference
of their squared norms times \(dV_h\), namely
\((|W^\pm|^2+R_h^2/48-|T_h|^2/4)dV_h\).
Subtracting the boundary transgression gives the relative second
Chern integer, so
\[
 I_\pm(U,h)=s_\pm\bigl(P_0k_\pm+C_\pm(\partial U,h)\bigr),
 \qquad k_\pm\in\mathbb Z .
\]
The transgression identity used here is
\cite[\S3, Proposition 3.2]{ChernSimons1974Characteristic}.
Using the canonical conformal spin isometry to transport the frames
preserves \(k_\pm\). Hence
\begin{equation}\label{alt:gauge-charge}
 I_\pm(U,h)-I_\pm(U,g)=-2J,\qquad
 C_\pm(\partial U,h)-C_\pm(\partial U,g)=-2s_\pm J .
\end{equation}
Independent frames give the second equation modulo \(P_0\).
For a single face the real equality follows by applying the
formula to a collar with the conformal change cut off at its
other end. In particular, using \(d_P(x)\) for the distance
to \(P\mathbb Z\), one gets
\[
 d_{P_0}\bigl(C_\pm(\partial U,h)\bigr)
 \le d_{P_0}\bigl(C_\pm(\partial U,g)\bigr)+2d_{P_0/2}(J).
\]
For example, if \(p_h\ge0\) on \(U\), then
\(0\le E_\pm(h;U)\le |I_\pm(U,g)|+2\max\{-J,0\}\).
In a volume-one Yamabe source the positive part of \(J\) is
bounded by \(Y^2\operatorname{Vol}_g(U)/96\) by \eqref{alt:schouten}.
Controlling the other sign is an additional boundary estimate.

Global curvature and volume bounds alone leave that sign open
on an actual residual. On the unit round sphere let
\(h_\lambda=\Phi_\lambda^*g\), where in stereographic coordinates
\(\Phi_\lambda(x)=\lambda x\), and take
\(U_\lambda=\{a/\lambda\le |x|\le b/\lambda\}\) for fixed
\(0<a<b\). The two gauges have the same global curvature bounds
and positive constant \(p=3/2\). The physical \(g\)-volume of
\(U_\lambda\) tends to zero and its \(h_\lambda\)-volume is the
fixed round volume of \(\{a\le |x|\le b\}\). Thus
\(J=q_g(U_\lambda)-q_{h_\lambda}(U_\lambda)\) has a negative
nonzero limit. This example concerns the zero Weyl class.

\begin{proposition}[Positive chiral charge on a residual domain]
\label{alt:positive-charge}
If \(p_h>0\) globally, put \(U^c=\overline{M\setminus U}\).
Then
\[
 d_{P_0}\bigl(C_\pm(\partial U,h)\bigr)
 =\min\{I_\pm(U,h),I_\pm(U^c,h)\}
 \ge\min\{E_\pm(h;U),E_\pm(h;U^c)\}.
\]
In particular, along a sequence with complementary chiral mass
bounded below and modular distance tending to zero, the chiral
Weyl energy and the Schouten integral on \(U\) both tend to zero.
\end{proposition}

\begin{proof}
The two charges are nonnegative and sum to \(P_0\) by
\eqref{alt:schouten}. The transgression congruence gives the
minimum exactly. If \(I_\pm(U^c,h)\) is bounded below by a positive
constant, the smaller charge is eventually \(I_\pm(U,h)\).
In the constant-\(p_h\) gauge,
\[
 I_\pm(U,h)=E_\pm(h;U)+(P_0-E)
                \frac{\operatorname{Vol}_h(U)}{\operatorname{Vol}_h(M)}.
\]
Thus the volume fraction is bounded by the modular distance divided
by \(\delta\) once the smaller charge is selected.
The same argument works in an iterated exhaustion
\(\lim_{L\to\infty}\limsup_i\), provided the complementary lower
bound holds in \(\liminf_i\) with one constant for all sufficiently
large fixed \(L\).
\end{proof}

A single balanced cut has modular distance at least the smaller of
its two Weyl masses. For a middle residual one needs the oriented
sum on its actual boundary. On a radial spherical collar
\(h=e^{2w}(dt^2+\gamma_{S^3(1)})\), let \(v=w'\).
The radial and tangential sectional curvatures are
\(-e^{-2w}v'\) and \(e^{-2w}(1-v^2)\). Since \(W=0\) and
\(p_h=3k_rk_t/2\), integration gives
\[
 I_\pm([t_-,t_+]\times S^3,h)=F(v(t_+))-F(v(t_-)),\qquad
 F(v)=4\pi^2-6\pi^2v+2\pi^2v^3.
\]
For the slice with increasing-\(t\) coorientation the modular
Chern--Simons value is \(s_\pm F(v)\). To fix the additive
constant, fill the slice from below by a radial ball whose metric
agrees near the slice and has \(w=t+\mathrm{const}\) at the
early pole. That pole has slope one and \(F(1)=0\).
For \(|v|\le1\) the modular distance on the slice is
\(2\pi^2(1-|v|)^2(2+|v|)\). Each boundary of an exact round
product cylinder has distance \(4\pi^2\), while the oriented
sum is zero. Strictly admissible round-class examples with
the same limiting cancellation are supplied by
\[
 w_L(t)=L-\log 2-\tfrac12\bigl(\log\cosh(t+L)+\log\cosh(t-L)\bigr).
\]
Here \(-1<v_L<1\), \(v_L'<0\), both pole expansions are smooth,
and \(w_L,v_L\to0\) uniformly on \([-L/2,L/2]\). The round
strip \([-L,0]\) with \(w=-\log\cosh t\) instead has limiting
charge \(4\pi^2\). The symmetric strip \([-L,L]\) in that
metric has charge tending to \(P_0\), so its modular distance
vanishes on the complementary branch. These examples explain
the paired orientation and the uniform complementary mass
hypothesis of Proposition~\ref{alt:positive-charge}.
For example, in a fixed admissible perturbation of the round
metric with nonzero Weyl tensor, take \(U\) to be the exterior
of a ball shrinking to a point with nonzero curvature in the
tested block. The complementary mass is positive and tends to
zero, while the exterior Weyl mass tends to the positive mass
of the whole sphere.

Even the local equation \(p_h=q>0\) permits two large individual
fluxes with a small oriented sum. On \([-L,L]\) set
\(w(0)=-\log L\), \(v(0)=0\) in
\[
 w'=v,\qquad v'=-\tfrac{2q}{3}e^{4w}/(1-v^2).
\]
As long as \(|w+\log L|\le1\) and \(|v|\le1/2\),
\(|v'|\le C_q L^{-4}\). Integration closes these bounds for
large \(L\) and gives \(|v|\le C_qL^{-3}\) and
\(|w+\log L|\le C_qL^{-2}\) on the full interval.
The individual modular distances tend to \(4\pi^2\), and their
paired charge is \(O_q(L^{-3})\). In this solution the first
integral \((1-v^2)^2-\tfrac{2q}{3}e^{4w}\) tends to one as
\(L\to\infty\). A smooth pole in a global radial completion
with constant \(q\) would force that first integral to be zero.

\section{Round-base restart and the Green-rescaling budget}
\label{alt:base}

Let \(I^+(M)\) be the infimum of \(E_+\) over smooth
positive-Yamabe classes on a fixed oriented smooth homotopy
four-sphere \(M\). Consider volume-one Yamabe minimizers \(g_i\)
with common chiral energy \(E_i\in[4\pi^2,8\pi^2-\delta]\),
\(E_i\to e\), \(Y_i\to y>0\), and
\(\mathfrak B(g_i):=\int_M|B_{g_i}|\,dV_{g_i}\to0\).
We first use a physical comparison with a round base, then a
comparison in a potentially different conformal gauge.

\subsection{A restart at one physical puncture}

\begin{theorem}[Restart on the original smooth type]\label{alt:restart}
Suppose \(M\) is nonstandard, \(e=I^+(M)\) and
\(y=Y_{\mathrm{rd}}\). Supply a volume-one round sphere
\((N,q_N)\) with a finite marked set \(P\), and oriented embeddings
on neighborhoods of whole fixed compact truncations of
\(N\setminus P\) on which the physical metrics converge smoothly
to \(q_N\). Then there is a slow coherent exhaustion, and after
a subsequence one fixed puncture label, whose actual exterior
contains \(e-o(1)\) of both chiral Weyl masses. Its intrinsic
physical diameter tends to zero. Capping that exterior and
selecting a volume-one Yamabe minimizer gives metrics \(g_i'\)
on the original oriented smooth type with
\[
 E_+(g_i')\to I^+(M),\qquad
 \mathfrak B(g_i')\to0,\qquad Y([g_i'])\to Y_{\mathrm{rd}} .
\]
Each restarted class has a marked conformally round germ.
\end{theorem}

\begin{proof}
A comparison on the whole unpunctured sphere would imply
\(E_i\to0\), so \(P\ne\varnothing\). Use disjoint round normal
coordinate balls with radius variable \(\rho\) around its points.
At cut radius \(r\), set \(t=\log(r/\rho)\) and choose the angular
orientation consistently. In the \(\rho^{-2}\) gauge the round
comparison is
\[
 dt^2+\left(\frac{\sin(k\rho)}{k\rho}\right)^2\gamma_{S^3(1)},
\]
where \(k^2>0\) is the round sectional curvature. The product
\(C^2\) error through \(|t|\le\ell+2\) is
\(O((r e^{\ell+2})^2)\). The \(C^4\) error on the central unit
splices is \(O(r^2)\). At fixed \(r,\ell\) the smooth physical
comparison gives an additional \(o_i(1)\) in both estimates.

Choose lengths \(\ell_m\to\infty\), tolerances
\(\eta_m\downarrow0\) with \(\ell_m\eta_m^2\to0\), then radii
\(r_m\downarrow0\) small enough for the buffered closures to be
disjoint and the round errors below \(\eta_m/2\). Also require
\(r_m^4(1+\ell_m)\to0\). Choose increasing source thresholds at which
the errors of the source on the entire larger truncation
\(N\setminus\bigcup_{p\in P} B_{q_N}(p,r_m e^{-\ell_m-2})\)
are below \(\eta_m\) and the volume integral over the truncation
at \(r_m\) agrees with its round value to within \(\eta_m\).
Impose unscaled smooth comparison through order \(m+4\) at each
threshold with error tending to zero. At each index use only the
map for the latest threshold and its restrictions. Write
\(r_i',\ell_i'\) for the resulting cut radius and buffer length.
This choice supplies the nested truncations and their comparisons
on all fixed compact subsets.

By the whole-truncation topology in the proof of
Theorem~\ref{thm:finite-punctures}, there is one exterior \(V_{ij}\)
per face and its marked spherical cap is a homotopy sphere.
Capping the central truncation gives \(N\), and the marked
connected sum of all capped pieces gives \(M\).
Let \(m^\pm_{ij}=E_\pm(g_i;V_{ij})\) and let \(e^\pm_{ij}\)
be the capped chiral energies. The buffered cap estimates in
Theorem~\ref{thm:caps} give
\begin{equation}\label{alt:restart-cap}
 \begin{split}
 |e^\pm_{ij}-m^\pm_{ij}|&\le C\eta_{2,i}^2,\\
 \mathfrak B(\widehat g_{ij})&\le
         \mathfrak B(g_i;V_{ij})+C\eta_{4,i},\\
 Y([\widehat g_{ij}])&\ge
   \min\{Y_i,(1-C\eta_{2,i})Y_{\mathrm{rd}}\}
        -C(\ell_i')^{-3/2}.
 \end{split}
\end{equation}
Here \(\eta_{2,i},\eta_{4,i}\le\eta_m\) along the chosen
exhaustion. The first bound compares both old and new unit splices
to the conformally flat product. The second uses the fourth jets
on the central splice and the conformal invariance of the Bach
\(L^1\) density in dimension four. Both bounds are independent of
\(r_i'\). One may put the splice and cap at physical length
scale \(r_i'\), using the factor \(r_i'/\rho\) on the retained
collar, extended with a cutoff on the spare buffer. The cap in this
pre-minimization representative has volume \(O((r_i')^4)\)
and the replaced representative on the long collar has volume
\(O((r_i')^4(1+\ell_i'))\).

All capped exteriors have positive Yamabe constant with limits
\(Y_{\mathrm{rd}}\), by \eqref{alt:restart-cap} and the universal
upper bound. The central capped piece carries the round metric.
By Theorem~\ref{thm:nonsplitting}, the strict upper band and the
\(o(1)\) cap costs imply at most one capped piece of energy
\(\ge4\pi^2\). The nonstandard type forces one nonstandard
capped exterior, of the original oriented type. For its label
\(j_i\),
\[
 I^+(M)\le e^+_{i,j_i},\qquad
 E_i-m^+_{i,j_i}\le E_i-I^+(M)+C\eta_{2,i}^2\to0 .
\]
The minus estimate follows because both closed energies are equal
on each capped piece and on each source. Fix the surviving label
along a subsequence. Larger cuts in the same normal coordinate
range, defined by restriction of the same map, add a marked
cylindrical layer. The capped type is preserved and the larger
exterior still contains the limiting mass. Smaller cuts within
the controlled buffer lose at most
\(C(1+\ell_i')\eta_{2,i}^2=o(1)\) of either Weyl mass.

The exteriors have physical volume tending to zero by the
volume-one normalization and the imposed convergence of the
central volume. The source Yamabe inequality gives uniform
physical lower bounds \(\operatorname{Vol} B(x,u)\ge c_0u^4\)
for \(0<u\le r_0\). In fact the cutoff in a ball and the
bounded source scalar curvature give \(V_x(u)/u^4\ge
c(V_x(u/2)/(u/2)^4)^{1/2}\). Iteration to small radii gives
the uniform bound. For each exterior point, a ball of radius
\(\min\{d(x,\partial V_{ij})/2,r_0\}\) lies in its exterior.
Its small volume thus implies
\(d(x,\partial V_{ij})\le 2c_0^{-1/4}
\operatorname{Vol}(V_{ij})^{1/4}\) for large \(i\).
Shortest segments to the boundary stay in the exterior up to the
first boundary point. Its sole boundary sphere has intrinsic
diameter \(O(r_i')\), giving the intrinsic diameter bound
\[
 \operatorname{diam}_{\rm intr} V_{ij}
 \le 4c_0^{-1/4}\operatorname{Vol}(V_{ij})^{1/4}+Cr_i'\to0 .
\]

Transport the surviving capped class to \(M\) by the oriented
diffeomorphism and take a unit-volume Yamabe minimizer
\cite[\S3]{Schoen1984Yamabe}. The class invariants are unchanged.
Equation \eqref{alt:restart-cap} and the minimizing bound imply the
three asserted limits, and the exact cap supplies the marked germ.
\end{proof}

For a given \(i\) the signed change equals the splice excess
\(e^+_{i,j_i}-m^+_{i,j_i}\) minus the discarded mass
\(E_i-m^+_{i,j_i}\). A strict improvement requires the second
term to exceed the first. A further restart requires new usable
comparison data. An estimate for the size of the marked germ
in the restarted Yamabe gauge requires further control.

\subsection{A conformal Green model}

For the same source hypotheses preceding
Theorem~\ref{alt:restart}, allow general \(y>0\) and \(e\).
Supply a smooth closed oriented homotopy-sphere conformal base
\((N,h)\) with finite marked set \(P\). Suppose on full compact
truncations off \(P\) there are oriented comparison embeddings
\(F_i\) and representatives \(F_i^*g_i=u_i^2h_i\), where
\(h_i\to h\) smoothly on fixed compacts. Compact-support Yamabe
tests then give \(Y([h])\ge y\) by cutoff at the finitely many
smooth points. The invariance of the Bach \(L^1\) density gives
\(B_h=0\) away from \(P\), and hence everywhere by smoothness.
Write \(e_N=E_\pm(h)\le e\), by conformal measure invariance
and exhaustion.

\begin{theorem}[One-pole Green budget]\label{alt:green}
Suppose a further constant length normalization at scales
\(\lambda_i\downarrow0\) gives
\[
 \lambda_i^{-2}F_i^*g_i\longrightarrow g_X=w^2h,\qquad
 w=\sum_{j=1}^{\ell}a_jG_{p_j},\quad a_j>0,
\]
smoothly on fixed compact subsets of \(N\setminus P\).
Here \(p_j\in P\) are distinct poles and \(G_p\) is the Green
function satisfying
\((6\nabla^*\nabla+R_h)G_p=24\pi^2\delta_p\).
On \(X=N\setminus\{p_1,\ldots,p_\ell\}\), including any ordinary
marked points, one has
\[
 \int_X |\operatorname{Ric}_{g_X}|^2\,dV_{g_X}
       =4e_N+32\pi^2(\ell-1),\qquad
 e-e_N\ge 8\pi^2\ell-y^2/48 .
\]
There is exactly one pole in the strict upper band, and hence
\(e-e_N\ge(Y_{\mathrm{rd}}^2-y^2)/48\).
\end{theorem}

\begin{proof}
A positive homogeneous solution with no pole is impossible on the
closed base with \(Y([h])>0\). The conformal scalar-curvature law
gives \(R_{g_X}=0\) away from the poles. Near each pole,
the corresponding Green function has positive \(r^{-2}\) leading
term in conformal normal coordinates
\cite[Lemma 6.4 and Theorem 6.5]{LeeParker1987Yamabe}. The other
terms are smooth. Inversion and a constant coordinate rescaling
give one Euclidean, unquotiented end per pole with differentiable
decay \(O_2(R^{-\alpha})\) for any fixed \(0<\alpha<2\).
The metric is complete and its curvature square is integrable
by the decay. A large end truncation removes \(\ell\) balls from
\(N\), so has Euler characteristic \(2-\ell\).
The Gauss--Bonnet boundary integral tends to \(8\pi^2\) on each
end: the rescaled end collar tends in \(C^2\) to the Euclidean
unit sphere collar with outward normal. The sign and coefficient
are those of a Euclidean ball
\cite[\S7, equation (7.2)]{TianViaclovsky2005ALE}. Since the
total chiral Weyl energies on \(X\) equal \(e_N\), the open
Gauss--Bonnet identity is
\[
 8\pi^2(2-\ell)=2e_N-\tfrac12
       \int_X|\operatorname{Ric}_{g_X}|^2\,dV_{g_X}
       +8\pi^2\ell .
\]

On an embedded fixed compact away from all of \(P\), the limit
Ricci integral equals the limit of physical trace-free Ricci
integrals, using only constant-scale invariance and smooth
convergence to a scalar-flat metric. Exhaustion proves the bound
by the limiting source budget
\[
 \lim_i\int_M|\operatorname{Ric}_{0,g_i}|^2,dV_{g_i}
     =4e+y^2/12-32\pi^2
\]
from \eqref{alt:schouten}. The ordinary omitted marks carry
no smooth curvature mass. The two Ricci identities give the
asserted inequality. Two poles would force \(e\ge e_N+8\pi^2\)
since \(y\le Y_{\mathrm{rd}}\), contradicting the strict band.
\end{proof}

Here the Green model and the compact conformal base have the same
Weyl measures on the comparison domains. A separate disjoint bubble
would need additional comparisons. If \([h]\) is nonround,
Theorem~\ref{thm:closed-gap} gives \(e_N>4\pi^2\) and
\(y\le Y([h])<Y_{\mathrm{rd}}\) by strict Yamabe rigidity
\cite[Theorem 2]{Schoen1984Yamabe}. Since
\(e-e_N<4\pi^2-\delta\), Theorem~\ref{thm:finite-punctures}
gives the same smooth type as \(M\) from the conformal comparison.
Thus \(I^+(M)\le e_N\), and the Green model is excluded at the
infimum. More generally the uniform exact-class Yamabe deficit
\(Y([h])\le Y_{\mathrm{rd}}-d_\delta\) in this closed band
gives
\[
 e-I^+(M)\ge e-e_N\ge
 \frac{Y_{\mathrm{rd}}^2-(Y_{\mathrm{rd}}-d_\delta)^2}{48}>0 .
\]
This uses the deficit on the closed base in
Lemma~\ref{thm:exact-compactness} and
\cite[\S12.2]{Baker2026TwistorSections}.

For a nonround base, the supplied zero-factor comparisons can be
normalized by the factor at an ordinary point. Harnack bounds and
interior regularity give a positive limiting factor solving the
homogeneous conformal Laplace equation away from \(P\). The
nonnegative isolated point sources give the Green representation
in Theorem~\ref{alt:green} under the smooth coherent comparison
hypotheses \cite{Baker2026MarkedLimits}. A positive unrescaled
limit factor instead extends through the marked points. Its
volume \(v\le1\) satisfies \(Y([h])v^{1/2}\le yv\) by the
Yamabe inequality and the limiting equation, so \(v=1\) and
\(Y([h])=y\). The physical metrics converge in the measured
Gromov--Hausdorff sense along that subsequence by whole-truncation
comparison, collapse of the one-face complements and source
ball noncollapse \cite{Baker2026MarkedLimits}.
In this physical branch at the infimum all chiral,
scalar-square, trace-free Ricci and full curvature defects vanish
by volume and energy saturation and \eqref{alt:full-curvature}.
For a round base, if the Green comparison is supplied then
\(e_N=0\) and the single-pole Ricci identity makes the complete
scalar-flat conformally flat model Euclidean. Other ordinary
marks can still be omitted by the comparison. The resulting
flat model carries zero chiral mass and the source may retain
energy outside the comparison.

\section{Closed critical metrics}
\label{alt:closed}

\subsection{Isolation of the first exact critical value}

For a fixed smooth closed oriented homotopy four-sphere \(M\), let
\[
 \mathcal A_M=\{E_+([g]): B_g=0,\ Y([g])>0,\ 0<E_+([g])<8\pi^2\}
\]
range over smooth metrics on \(M\), and set
\(a_M=\inf\mathcal A_M\) when the set is nonempty. Let
\(\kappa_\delta>0\) be the exact same-type energy drop in the
upper band \(E_+\le 8\pi^2-\delta\) from
Theorem~\ref{thm:ale-descent}.

\begin{theorem}[First fixed-type critical value]\label{alt:critical}
If \(\mathcal A_M\ne\varnothing\), then its first value is attained
and \(4\pi^2<a_M<8\pi^2\). For
\[
 a_M\le b\le8\pi^2-\delta,\qquad b<a_M+\kappa_\delta ,
\]
the volume-one Yamabe representatives realizing values in
\(\mathcal A_M\cap[a_M,b]\) are sequentially compact up to oriented
diffeomorphisms. The set of distinct values in that sublevel is
finite. In particular \(a_M\) is isolated among exact positive
subcritical values on \(M\).
\end{theorem}

\begin{proof}
Take a sequence of exact values converging to \(e\le b\).
Theorem~\ref{thm:closed-gap} places the representatives in the
exact compactness band. In the blow-up alternative of
Lemma~\ref{thm:exact-compactness}, the smooth ALE compactification
and the marked same-type descent produce a member of
\(\mathcal A_M\) of value at most
\(e-\kappa_\delta<a_M\), which is impossible. The compact
alternative gives smooth convergence with limiting energy \(e\)
and positive Yamabe sign. Apply this first to a minimizing
sequence, taking \(\delta>0\) with \(a_M<8\pi^2-\delta\) and
discarding a finite prefix, to attain \(a_M\). The closed gap
gives the strict lower bound at the attaining class.

For finiteness, fix any smooth Bach-flat metric \(g\) on \(M\)
and take the trace-free Bach gauge of Gursky--Viaclovsky
\cite[\S2.3, equation (2.59) and Theorem 2.11]{GurskyViaclovsky2015}.
For \(0<\alpha<1\) its map is
\[
 P_g^B(\theta)=\pi_0^g\!
     \left(B_{g+\theta}
         +K_{g+\theta}(\square_{K_g}\beta_g\theta)\right)
 : C^{4,\alpha}(S^2_0)\longrightarrow C^{0,\alpha}(S^2_0)
\]
with the normalization of that source. Here \(K_g\) is its conformal
Killing operator, \(\square_{K_g}=\delta_g K_g\),
\(\beta_g\theta=\delta_g\theta\) on trace-free tensors, and \(\pi_0^g\)
is the fixed-background trace-free projection. This map is real
analytic in \(\theta\): inverse metric, differentiation with fixed
background, and multiplication through the required jet orders
are analytic on the open cone of positive metrics. Theorem 2.11
gives an elliptic derivative and puts every sufficiently
\(C^{4,\alpha}\)-close Bach-flat metric into the small zero set by
a diffeomorphism near the identity and a conformal change.
The elliptic derivative has a finite-dimensional kernel \(K\)
and closed image of finite codimension. Project onto that image
and apply the analytic implicit function theorem on a closed
complement of \(K\). Its graph \(\theta=k+v(k)\) contains all
small zeros.

The full four-index Weyl integral \(F(k)\) on the graph is
real analytic by the same metric operations and the analytic
volume density. At a parameter from an actual Bach-flat metric
its unrestricted first variation vanishes, so \(dF(k)=0\) on
the entire kernel tangent space. The finite-dimensional
\L{}ojasiewicz gradient inequality
\cite[Theorem 1]{Feehan2020Gradient} gives
\(|F(k)-F(0)|^\vartheta\le C|dF(k)|\), with
\(0<\vartheta<1\) near zero. Consequently all nearby smooth
Bach-flat metrics have the same full Weyl integral, and hence
the same \(E_+\) on \(M\). If the displayed sublevel contained
infinitely many distinct values, choose a representative
of each and use smooth subsequential compactness. Local constancy
at its limit contradicts distinctness.
\end{proof}

The isolation concerns values. A least exact critical value
may be a saddle for variations through arbitrary positive-Yamabe
metrics. More generally, distinct exact values accumulating
inside a fixed band force the blow-up alternative along the
corresponding sequence, and an exact value lower by at least
\(\kappa_\delta\) than the accumulation value on the same type.

\subsection{The chiral Cotton obstruction and Einstein reduction}

Let \(C_{ijkl}=(W^+)_{ijkl}\) be the all-lowered Weyl block,
\(D_{jkl}=\nabla^i C_{ijkl}\), and
\(L(\alpha)_{jkl}=\alpha^iC_{ijkl}\) for a real one-form.
Put \(w=|W^+|_{\mathrm{HS}}\) and \(U=\{w>0\}\).
Use the full ordered-index norm for three-tensors.
In an orthonormal self-dual eigenbasis \(\Sigma_a\), with
eigenvalues \(\lambda_a\),
\(C_{ijkl}=\sum_a\lambda_a(\Sigma_a)_{ij}(\Sigma_a)_{kl}\).
Contraction of each normalized eigenform with a vector has
squared norm half the vector's squared norm, and summing
over the ordered last pair contributes a factor of two.
Thus \(L^*L=w^2\operatorname{Id}\). Define on \(U\)
\[
 \beta=-w^{-2}L^*D,\qquad
 \mathcal R=D+L(\beta),\qquad r=|\mathcal R|,
\]
and set \(\mathcal R=D\) on \(M\setminus U\).
The conformal change \(\widehat g=e^{2f}g\) gives
\(\widehat C=e^{2f}C\) and \(\widehat D=D+L(df)\) as covariant
tensors \cite[equations (3.3)--(3.4)]{GurskyMcKeownTyrrell2026}.
Because \(\widehat L=L\), \(\widehat L^*=e^{-4f}L^*\), and
\(\widehat w^2=e^{-4f}w^2\), the residual \(\mathcal R\) is
unchanged and \(\widehat\beta=\beta-df\). The pointwise orthogonal
decomposition is
\begin{equation}\label{alt:cotton-split}
 |D+L(df)|_g^2=r^2+w^2|df-\beta|_g^2 \quad\text{on } U.
\end{equation}
The left-hand side is \(r^2\) on the complement.
In particular
\[
 \int_M |\widehat D|_{\widehat g}^{4/3}\,dV_{\widehat g}
 \ge\int_M r^{4/3}\,dV_g ,
\]
using the three-index norm's conformal weight. On any piecewise
smooth regular closed curve \(\gamma\subset U\), the integral
of the left side of \eqref{alt:cotton-split} with respect to
arclength is bounded below by
\[
 \int_\gamma r^2\,ds_g+
         \frac{|\oint_\gamma\beta|^2}{\int_\gamma w^{-2}\,ds_g}.
\]
This follows from the tangential component of \(df-\beta\),
Cauchy--Schwarz and \(\oint_\gamma df=0\).

For a smooth half-harmonic representative, \(\mathcal R=0\)
on \(U\) and \(\beta\) must extend to a smooth closed one-form.
Conversely such an extension is exact on a homotopy sphere.
The resulting divergence equation extends by continuity to
\(\overline U\), and holds on the interior of the complement
because \(C\) vanishes there.

\begin{theorem}[Closed harmonic scale for a Bach-flat class]
\label{alt:harmonic-scale}
Let \(g\) be smooth and Bach-flat on a closed oriented homotopy
four-sphere and suppose \(W^+\not\equiv0\). Then a smooth
half-harmonic representative for \(W^+\) exists if and only if
\(\mathcal R=0\) on \(U\) and \(\beta\) has a smooth extension
to the whole sphere. Every such half-harmonic representative
is Einstein.
\end{theorem}

\begin{proof}
With the Schouten tensor \(P_g\), the Bach formula in the
\(W^+\) block is
\(B_{ij}=2(\nabla^k\nabla^l+P_g^{kl})C_{ikjl}\)
\cite[\S3, p.~78]{LeBrun2013EinsteinExtremal}.
Use \(\nabla^l C_{ikjl}=-D_{jik}\) and substitute
\(D=\mathcal R-L(\beta)\) twice to obtain
\begin{equation}\label{alt:cotton-bach}
 \tfrac12 B_{ij}
 =C_{ikjl}\bigl(P_g^{kl}-\nabla^k\beta^l+\beta^k\beta^l\bigr)
       -\nabla^k\mathcal R_{jik}-\beta^k\mathcal R_{ikj}.
\end{equation}
Here the derivative of the Weyl tensor in the substitution
uses \(\nabla^k C_{ajik}=-D_{iaj}\).
When \(\mathcal R=0\), the skew part yields
\(W^+((d\beta)^+)=0\) by the Bianchi identity.

The invertible locus of \(W^+\) is dense in \(U\) when
\(D=-L(\beta)\) there. Indeed, if rank two held on an open
subset, a local quaternionic frame \(I,J,K=IJ\) could be
chosen with \(J\) the zero eigenline, and with opposite
eigenvalues on \(I,K\). Differentiate \(W^+(\omega_J)=0\)
and contract on the first index of \(C\). The contracted
derivative of \(C\) vanishes against \(\omega_J\) by
\(D=-L(\beta)\). Writing \(\nabla J=a\otimes I+b\otimes K\)
in the other term gives \(a(IX)=b(KX)\), hence \(b=a\circ J\).
The identity \((\nabla_{JX}J)JY=(\nabla_XJ)Y\) follows and
gives integrability of \(J\). For a Hermitian metric,
the \(W^+\) block consists of a scalar multiple of
\(\tfrac12\omega_J\otimes\omega_J-\tfrac13\operatorname{Id}\)
and off-diagonal terms between \(\omega_J\) and its
orthogonal complement. Thus \(W^+(\omega_J)=0\) makes the
entire block zero, contrary to rank two
\cite[\S4, Theorem 4.1 and proof of Proposition 4.1]{GoverNagy2007FourDimensional}.

Consequently \((d\beta)^+=0\) on \(U\).
For a nonzero smooth Bach-flat block, \(U\) is dense in \(M\):
the conformally normalized metric is real analytic in local
conformal harmonic coordinates by obstruction-flatness
\cite[Theorems 2.2 and 3.7(c)]{LassasLiimatainen2019ConformalHarmonic}.
The analytic identity theorem for a Weyl block and the conformal
invariance of its zero set exclude an open zero set on the
connected manifold. If \(\bar\beta\) is a smooth extension,
\(d\bar\beta\) is therefore anti-self-dual everywhere. Stokes
gives \(0=\int_M d\bar\beta\wedge d\bar\beta
=-\|d\bar\beta\|_2^2\). The closed extension is exact, giving
the asserted half-harmonic representative by \eqref{alt:cotton-split}.
The necessity follows by taking \(\bar\beta=df\).

In the half-harmonic scale the Bach equation is
\(C_{ikjl}P_g^{kl}=0\). On trace-free symmetric tensors this
contraction is \(W^+\) acting on the first factor in
\(S^2_0\simeq\Lambda^+\otimes\Lambda^-\).
Indeed let \(I_a,J_b\) be opposite-chirality unit quaternionic
frames with \(W^+\) diagonal in the former. The contraction on
a basis element \(I_aJ_b\) is
\(\frac12\sum_c\lambda_c I_c(I_aJ_b)I_c^{-1}\).
The conjugation fixes the term with \(c=a\) and negates the
others, yielding \(\lambda_a I_aJ_b\). Thus the trace-free
Schouten tensor is zero on the invertible locus, and then
everywhere by density. The contracted Bianchi identity gives
constant scalar curvature.
\end{proof}

The Einstein implication in the proof requires a nonzero
half-harmonic block and Bach-flatness, and applies on any connected
closed oriented four-manifold. It shows a limitation of using
stability to obtain a harmonic scale on general smooth types.
On \(N=\#_4\mathbb{CP}^2\) in the positive-signature orientation,
LeBrun supplies a positive-Yamabe self-dual conformal class
\cite[Theorem 1 and \S5, p.~235]{LeBrun1991Explicit}.
Its \(W^-=0\) makes \(E_+=12\pi^2\tau(N)=48\pi^2\)
an absolute minimum by the signature formula. It is Bach-flat
with nonnegative unrestricted Hessian. If its nonzero \(W^+\)
had a half-harmonic representative, the preceding implication
would make it positive Einstein. But then
\(4\pi^2(2\chi-3\tau)=2E_-+\int R^2/24\,dV>0\),
whereas \(\chi=6\), \(\tau=4\). On the homotopy sphere in the
subcritical band, both the vanishing residual and smooth extension
in Theorem~\ref{alt:harmonic-scale} remain hypotheses.

\subsection{The half-harmonic energy threshold}

\begin{theorem}[Strict half-harmonic bound]\label{alt:harmonic-gap}
Let \([g]\) be a smooth positive-Yamabe class on a closed
oriented homotopy four-sphere, with a smooth representative
satisfying \(\nabla^i(W^+)_{ijkl}=0\). Let \(e=E_+([g])\),
and let \(T_Y\) be the squared \(L^2\) norm of trace-free
Ricci curvature in a Yamabe minimizer of the class.
Then either the class is conformally round or
\[
 e>\frac{16\pi^2}{3}+\frac{T_Y}{6}.
\]
For the largest eigenvalue \(z\) of \(W^+\) in the nonround
case, the operator \(-6\Delta_g+R_g-6z\) has negative first
eigenvalue, and \(36\int_M z^2\,dV_g>Y([g])^2\).
\end{theorem}

\begin{proof}
In the half-harmonic representative put \(w=|W^+|_{\mathrm{HS}}\).
Use \(\Delta_g=\nabla^i\nabla_i\). The half-harmonic
Weitzenb\"ock formula and improved Kato inequality give
\[
 \tfrac12\Delta_g w^2=|\nabla W^+|^2+\tfrac12R_g w^2-18\det W^+,
 \qquad |\nabla W^+|^2\ge\tfrac53|\nabla w|^2\quad(w>0)
\]
\cite[equations (3.3), (3.7) and Corollary 3.4]{GurskyMcKeownTyrrell2026}.
Fu gives the gap with the four-index norm, whose square is
\(4w^2\) for one block
\cite[Proposition E, equations (2-1)--(2-9)]{Fu2018FourManifolds}.
For eigenvalues \(x\le y\le z\) with \(x+y+z=0\),
\[
 6zw^2-36\det W^+=12z(x-y)^2\ge0,\qquad
 6z\le 2\sqrt6\,w.
\]
The first equality uses \(w^2=2z^2-2xy\). The bound uses
\(x^2+y^2\ge z^2/2\).
Set \(u_\epsilon=w^2+\epsilon^2\),
\(f_\epsilon=u_\epsilon^{1/6}\), and \(P_z=-6\Delta_g+R_g-6z\).
Then
\[
 6\Delta_g f_\epsilon=u_\epsilon^{-5/6}
   \bigl(\Delta_g u_\epsilon-\tfrac56u_\epsilon^{-1}
                            |\nabla u_\epsilon|^2\bigr).
\]
Kato bounds the subtracted term by \(2|\nabla W^+|^2\), also
at zeroes since \(\nabla u_\epsilon=0\) there. The two preceding
displays give
\(P_z f_\epsilon\le(R_g-6z)\epsilon^2 u_\epsilon^{-5/6}
\le\|R_g\|_\infty\epsilon^{1/3}\).
Multiplication by \(f_\epsilon\) and integration bound
\(f_\epsilon\) in \(H^1\). Its uniform limit \(f=w^{1/3}\)
therefore has \(P_z\) quadratic form at most zero by lower
semicontinuity. If \(w\not\equiv0\), the Yamabe inequality
and Hölder yield
\[
 Y\|f\|_4^2
 \le \int_M(6|\nabla f|^2+R_g f^2)\,dV_g
 \le 6\int_M zf^2\,dV_g
 \le 2\sqrt6\,\|w\|_2\|f\|_4^2.
\]
In equality, Hölder forces \(w\) proportional to \(w^{2/3}\)
and hence a positive constant by continuity and connectedness.
The constant test is Yamabe minimizing, giving constant scalar
curvature \(R_g=2\sqrt6\,w\). The Weitzenb\"ock formula and
\(18\det W^+\le\sqrt6\,w^3\) force \(W^+\) to be parallel,
with eigenvalues \((-R_g/12,-R_g/12,R_g/6)\). The simple
eigenline is orientable since \(H^1(M;\mathbb Z/2)=0\).
A global unit section is a nonzero parallel self-dual form,
contrary to \(b_2(M)=0\). Thus \(e>Y^2/24\).
In the Yamabe gauge the topological identity is
\(16\pi^2=2e-T_Y/2+Y^2/24\). Combining gives the asserted
strict bound. If \(W^+=0\), the zero signature gives \(W^-=0\)
and the simply connected closed conformally flat manifold is
conformally round by the developing map.

For the spectral statement, the form tested on \(f\) gives
\(\lambda_1(P_z)\le0\). If equality holds then \(f\) is a
weak first eigenfunction. Elliptic regularity and the strong
maximum principle for the bounded Lipschitz potential give
\(f>0\). The calculation at \(\epsilon=0\) now gives
\[
 P_z f=-w^{-5/3}
     \bigl(2|\nabla W^+|^2-\tfrac{10}{3}|\nabla w|^2
                    +12z(x-y)^2\bigr).
\]
The nonnegative summands force \(x=y=-z/2\) everywhere.
A global section of norm \(\sqrt2\) of the simple eigenline
defines a compatible almost complex structure. Its first
Chern class must satisfy
\(\langle c_1^2,[M]\rangle=2\chi+3\tau=4\), impossible
when \(H^2(M;\mathbb Z)=0\). Thus the first eigenvalue is
negative. A negative-form test and Hölder give
\(Y^2<36\int z^2\,dV_g\). The sign of the operator is conformally
invariant because its quadratic form is that of the conformal
Laplacian with the weight-two potential \(6z\).
\end{proof}

Every positive Einstein metric satisfies the half-harmonic
hypothesis by the contracted Bianchi identity. Thus in the band
\(4\pi^2<E_+\le16\pi^2/3\) every Bach-flat candidate lacks
the smooth harmonic scale of Theorem~\ref{alt:harmonic-scale}.

\subsection{The positive Einstein Hessian}

\begin{proposition}[A spectral criterion for Weyl descent]
\label{alt:einstein-hessian}
On a closed oriented homotopy four-sphere let \(g\) be positive
Einstein with constant scalar \(R\). Use the positive-leading
Lichnerowicz operator on transverse-traceless tensors
\[
 \mathcal L=\nabla^*\nabla-2\mathring{\operatorname{Rm}}
                       +\tfrac R2\operatorname{Id},
 \qquad (\mathring{\operatorname{Rm}}h)_{ij}=R_{ipjq}h^{pq},
\]
where \(\mathring{\operatorname{Rm}}h=-h\) on trace-free tensors
at sectional curvature one. The Hessian of \(E_+\) on this
subspace is
\begin{equation}\label{alt:hessian}
 D^2E_+(g)[h,h]=\tfrac18
         \langle h,(\mathcal L-R/2)(\mathcal L-R/3)h\rangle_{L^2}.
\end{equation}
A negative variation modulo conformal and diffeomorphism
directions exists if and only if the transverse-traceless spectrum
of \(\mathcal L\) meets \((R/3,R/2)\). Such a direction
decreases the energy inside the positive-Yamabe classes on \(M\).
\end{proposition}

\begin{proof}
Chern--Gauss--Bonnet and signature give the functional identity
\[
 \int_M(|\operatorname{Ric}|^2-R^2/3)\,dV
 =\tfrac12\int_M |W|_{\mathrm{tens}}^2,dV-16\pi^2\chi(M)
 =4E_+-16\pi^2\chi(M).
\]
Gursky--Viaclovsky's transverse-traceless factorization for this
quadratic functional at an Einstein metric gives half the product
of the operators \(-\mathcal L+R/2\) and \(-\mathcal L+R/3\)
\cite[Theorem 3.10 and equation (3.41)]{GurskyViaclovsky2015}.
Dividing by four gives \eqref{alt:hessian}. On an eigentensor
the multiplier has negative sign exactly in the stated open
interval, and the bilinear kernel consists of the two endpoint
eigenspaces. The \(R/2\) endpoint is the infinitesimal Einstein
equation in transverse-traceless gauge.

Conformal and Lie-derivative directions are in the radical
of the Hessian by differentiating the invariances at the
critical metric. The remaining part of a smooth variation is
transverse-traceless: remove its trace and solve the orthogonal
projection onto the image of the conformal Killing operator.
Its normal operator is elliptic, and the right-hand side is
orthogonal to the conformal Killing fields. A negative
transverse-traceless tensor is realized by the path
\(g_s=g(e^{s h^\sharp}\cdot,\cdot)\), whose volume form is fixed.
For small \(|s|\), scalar curvature stays positive by compactness
and smooth dependence, giving positive Yamabe sign.
\end{proof}

Exclusion of both endpoint eigenvalues gives local Bach
rigidity, and exclusion of the full closed interval gives a
strict local minimum modulo the two gauges by the nonlinear
rigidity and local minimality results of
\cite[Theorem 1.6 and Proposition 6.4]{GurskyViaclovsky2015}.
If the open interval is empty but an endpoint occurs, the
Hessian is nonnegative and leaves those kernel directions
for higher-order analysis. The criterion concerns unrestricted
metric variations with positive Yamabe sign. It gives no path
of exact Bach-flat classes descending from an isolated value.

For any positive Einstein representative of volume \(V\),
\(R\sqrt V=Y([g])\). To see this, write
\(g=u^{-2}\widehat g\) with \(\widehat g\) a Yamabe minimizer.
The trace-free Ricci equation is
\(0=\operatorname{Ric}_{0,\widehat g}
      +2u^{-1}(\widehat\nabla^2u)_0\).
Pair with \(u\operatorname{Ric}_{0,\widehat g}\) and integrate
in \(\widehat g\). The Hessian term integrates to zero by the
contracted Bianchi identity and the constant scalar curvature.
Thus \(\widehat g\) is Einstein too. The identity
\(16\pi^2=2E_++R^2 V/24\) for the two Einstein metrics
gives the claim. In particular in the unit-volume gauge,
\(R=Y\), \(E_+=8\pi^2-Y^2/48\), and the interval test is
\(\mu/Y\in(1/3,1/2)\). For sectional curvature one,
\(R=12\) and \(\mathcal L=\nabla^*\nabla+8\), so the Hessian
is \((\nabla^*\nabla+2)(\nabla^*\nabla+4)/8\) and is
positive on this subspace. For a nonround positive Einstein
survivor the integral Weyl band gives no conclusion about an
eigenvalue in the descent interval.
